\documentclass[11pt]{article}
\usepackage[a4paper,margin=25mm]{geometry}
\usepackage{amsmath,amssymb,hyperref,lmodern}
\title{A capacity-one fixed-query propagation lower bound}
\author{Working research note}
\date{16 November 2026}
\begin{document}
\maketitle
\paragraph{Scope.} The adopted protocol is specified in
\texttt{generic-proof.tex}. Time counts independent uniform global
initiations. The query is conditioned on a nonempty cache. This result has
passed independent internal mathematical review, not formal verification
or external peer review.
\section{A fixed-query propagation lower bound at capacity one}
Consider $c=L=1$, $n\ge3$, and start with the directed zero-age cycle
$C_i=\{(i+1\bmod n)^0\}$ on labels $0,\ldots,n-1$. Let $\nu_{i,t}$ be the
one-identifier query law at observer $i$, conditional on a nonempty cache.
For every real $R>0$,
\[
 \|\nu_{i,t}-\operatorname{Unif}([n]\setminus\{i\})\|_{\rm TV}
 \ge\left[1-
 \frac{n(1+1/n)^t}{(n-1)R}-\frac{2R}{n-1}\right]_+.
\]
Optimizing over $R$ gives the explicit bound
\[
 \|\nu_{i,t}-\operatorname{Unif}([n]\setminus\{i\})\|_{\rm TV}
 \ge\left[1-\frac{\sqrt{8n}(1+1/n)^{t/2}}{n-1}\right]_+.
\]
Thus a time guaranteeing query error at most $\varepsilon$ uniformly
over starts, for $0<\varepsilon<1$, must satisfy
\[
 t\ge\frac{\log\!\left((1-\varepsilon)^2(n-1)^2/(8n)\right)}
 {\log(1+1/n)}
 \quad\text{when the logarithm is positive}.
\]
For fixed $\varepsilon$ this is $n\log n-O_\varepsilon(n)$.
In particular, for every fixed $0<\delta<1$, at
$t_n=\lfloor(1-\delta)n\log n\rfloor$ this error tends to one. Hence for
any fixed $0<\varepsilon<1$, a worst-start fixed-observer query mixing time
has lower limit at least one after division by $n\log n$. This is a
sample-level result, separate from the stronger known full-state depth
lower bound, and no capacity-two extension is asserted.

\paragraph{Metric-length drift.}
Use circular distance $d(i,j)=\min(|i-j|,n-|i-j|)$, and define
$W(x)=\sum_{i:C_i\ne\varnothing}d(i,f(i))$, where $f(i)$ is the sole cached
identifier. Initially $W=n$. If a nonempty active peer $p$ contacts
$q=f(p)$ and $r=f(q)$ exists with $r\ne p$, the capacity-one identity update
is $f'(q)=p$, $f'(p)=r$. The change of $W$ is
$d(p,r)-d(q,r)\le d(p,q)$ by the triangle inequality. If $q$ is empty,
$W$ is unchanged. If $r=p$, the self-reply is removed and $W$ decreases by
$d(p,q)$. Empty initiations are no-ops. Averaging over the uniform initiator
therefore gives
\[
 \mathbb E[W(X_{t+1})\mid X_t]\le(1+1/n)W(X_t),
 \qquad \mathbb EW(X_t)\le n(1+1/n)^t.
\]

\paragraph{From length to the conditional query.}
The initial law and transition kernel are invariant under cyclic label
rotations. Thus every observer has the same availability probability
$a_t$ and the same expected distance, assigning zero distance to an empty
row. Weak connectivity and capacity one imply that there are at least
$n-1$ occupied rows: there must be at least $n-1$ underlying edges, and
there is at most one descriptor per row. Consequently $a_t\ge(n-1)/n$.
The conditional mean query distance at every observer is therefore
\[
 \mathbb E[d(i,\text{query})\mid C_i\ne\varnothing]
 =\frac{\mathbb EW(X_t)}{n a_t}
 \le\frac{n(1+1/n)^t}{n-1}.
\]
Markov's inequality gives conditional mass at least one minus this last
quantity divided by $R$ to the radius-$R$ ball around $i$. The uniform peer
law assigns that ball mass at most $2R/(n-1)$. Testing this event proves
the finite bound. The ball cardinality bound holds for every real $R>0$:
it is at most $2\lfloor R\rfloor\le2R$ below the diameter, and at most
$n-1\le2R$ at or above it. Substituting
$R=\sqrt{n(1+1/n)^t/2}$ proves the optimized bound. Rearrangement gives the
necessary time, and at $t_n$ its subtracted term tends to zero.

\end{document}
